\documentclass{beamer}

\useoutertheme{infolines}

\title[Introduction to LaTeX]{Introduction to LaTeX}
\subtitle{A Document Preparation System}
\author[Aryan Saraiya]{Aryan Saraiya}
\institute[Technical Writing]{Technical Writing}
\date{}

\begin{document}
\begin{frame}
    \titlepage
\end{frame}

\begin{frame}{Introduction to LaTeX}
    \begin{itemize}
        \item LaTeX is a document preparation system.
        \item It is widely used for academic and technical documents.
        \item It helps create well-structured and professional documents.
        \item It is especially useful for writing mathematical equations.
    \end{itemize}
\end{frame}

% Slide 3: Structure and Formatting
\begin{frame}{Structure and Formatting}
    A LaTeX document has two main parts:

    \begin{itemize}
        \item Preamble: Contains the document class, packages, and settings.
        \item Document body: Contains the actual content.
    \end{itemize}

    Some useful packages are:

    \begin{itemize}
        \item geometry: Sets page margins.
        \item sectsty: Changes chapter and section heading styles.
        \item fancyhdr: Adds headers and footers.
    \end{itemize}
\end{frame}

% Slide 4: Conclusion
\begin{frame}{Conclusion}
    \begin{itemize}
        \item LaTeX is useful for preparing academic documents.
        \item It provides consistent formatting and document structure.
        \item Packages make it easy to customize documents.
    \end{itemize}

    \vspace{0.5cm}
    \centering
    Thank You!
\end{frame}

\end{document}
